%% Shared body of the Supplementary Information.
%% Included by paper/supplementary.tex (npj Materials Degradation) and
%% paper-corrosion-science/supplementary.tex (Corrosion Science).  The two
%% wrappers differ only in the title; the content lives here so it cannot go
%% stale in one venue's copy.
\section{Model derivation}
\label{sec:derivation}

This section gives the complete derivation of the model used in the main text.
The mechanistic framework is the Point Defect Model (PDM) formulation of
Li et al.\ (2020) for oxide growth in supercritical water, specialized here from
supercritical to subcritical high-temperature water and applied to the duplex
oxide on Nb-stabilized X6CrNiNb18-10 (AISI~347) under the exposure conditions of
Veile et al.\ (2024). Equation numbers of the form ``Eq.~($n$) of the
supercritical-water formulation'' refer to that source; boxed (displayed and
labeled) equations are the working forms used in this study.

\subsection{System, geometry, and species}
\label{sec:geometry}

The scale is duplex. Coordinates are taken with the origin at the
barrier-layer/outer-layer (bl/ol) interface:
\begin{equation}
\underbrace{\text{metal (m)}}_{x = L_{\mathrm{bl}}}
\;\Big|\;
\underbrace{\text{barrier layer (bl), } \mathrm{MO}_{\chi/2}}_{x = 0}
\;\Big|\;
\underbrace{\text{outer layer (ol), } \mathrm{MO}_{\delta/2}}_{x = -L_{\mathrm{ol}}}
\;\Big|\;
\text{high-temperature water (e)} .
\end{equation}
That is, the metal/barrier-layer (m/bl) interface sits at $x = L_{\mathrm{bl}}$,
the bl/ol interface at $x = 0$, and the ol/water interface at
$x = -L_{\mathrm{ol}}$.

The microstructural identification of the PDM objects for X6CrNiNb18-10 in
boiling-water-reactor hydrothermal water, following the phase analysis of
Veile et al.\ (2024), is given in Table~\ref{tab:layers}.

\begin{table}[htbp]
\centering
\caption{Identification of the PDM layers with the observed oxide
microstructure on X6CrNiNb18-10.}
\label{tab:layers}
\resizebox{\textwidth}{!}{%
\begin{tabular}{llll}
\toprule
PDM object & Observed layer & Phase (XPS/EDX) & Valences \\
\midrule
barrier layer (bl) & inner nanocrystalline Cr-rich layer &
$\mathrm{FeCr_2O_4}$ / $\mathrm{NiCr_2O_4}$ spinel &
$\chi = 8/3$ (spinel-averaged cation) \\
outer layer (ol) & discrete Fe-rich crystals &
$\mathrm{Fe_3O_4}$ (magnetite), partly $\mathrm{Fe_2O_3}$ &
$\delta = 8/3$, $\theta = 3$ \\
--- (inert) & Ni enrichment zone at bl/matrix interface &
metallic/NiO traces & not modeled ($\mathrm{d}L_{\mathrm{Ni}}/\mathrm{d}t \approx 0$) \\
\bottomrule
\end{tabular}}
\end{table}

The Kr\"oger--Vink species in the barrier layer are: the cation vacancy
$V_M^{\chi\prime}$, the cation interstitial $M_i^{\chi+}$, the oxygen vacancy
$V_O^{\bullet\bullet}$, the oxide ion $O_O$, the lattice cation $M_M$, and the
electron $e'$. Constants: $\gamma = F/(RT)$; $\Omega_{\mathrm{bl}}$,
$\Omega_{\mathrm{ol}}$ are the molar volumes per cation of the barrier-layer
and outer-layer oxides; $\Omega_m$ is the molar volume of the metal; and
$\mathrm{PBR} = \Omega_{\mathrm{bl}}/\Omega_m$ is the Pilling--Bedworth ratio.

The environment corresponds to the conditions of Veile et al.\ (2024):
\begin{itemize}
\item $T = 513.15$~K (240~$^\circ$C), $P = 7$~MPa, i.e.\ liquid water with density
$\rho \approx 0.8137$~kg/L (IAPWS-97).
\item Since $\rho \gg 0.1$~g/cm$^3$, the electrochemical-oxidation
(high-density) branch of the supercritical-water PDM formulation applies: this
is the classical condensed-water PDM limit of the model.
\item Dissolved oxygen $[\mathrm{O_2}] = 0.4$~ppm; ultrapure water with
conductivity 0.055~$\mu$S/cm.
\item Neutral pH at temperature: $\mathrm{p}K_w(240\,^\circ\mathrm{C}) \approx 11.2$,
hence $\mathrm{pH} \approx 5.6$.
\item Dimensionless relative volumetric concentrations
(Eqs.~(13) and (14) of the supercritical-water formulation), with reference
concentration $C^0 = 1$~mol/L:
\begin{align}
C_{\mathrm{H_2O}} &= \frac{1000\,\rho}{18.015\, C^0_{\mathrm{H_2O}}} \approx 45.2 , \\
C_{\mathrm{O}} &= \frac{[\mathrm{O_2}]_{\mathrm{ppm}}\,\rho \times 10^{-3}}{32.00\, C^0_{\mathrm{O}}}
\approx 1.02 \times 10^{-5} .
\end{align}
\item Potential $V$: open circuit (no polarization applied in the autoclave).
$V$ is set by the $\mathrm{O_2/H_2O}$ redox couple (the electrochemical
corrosion potential) and is constant over an exposure; it is therefore absorbed
into effective rate constants (Section~\ref{sec:reduction}). The open-circuit
closure $i_{\mathrm{pol}} = 0$ (Eq.~(46) of the supercritical-water
formulation) becomes relevant only if $V$ is made explicit (parametric studies
in the electrochemical corrosion potential).
\end{itemize}

\subsection{Interfacial reactions}
\label{sec:reactions}

The complete reaction set of the supercritical-water PDM formulation is
retained, with the high-density/electrochemical-oxidation significance of each
reaction carried over to subcritical hydrothermal water. Here $\chi$ and
$\delta$ are the cation valences in the barrier and outer layers, and $\theta$
is the valence in the dissolved/further-oxidized product. The reactions at the
three interfaces are listed in Tables~\ref{tab:rxn-mbl}--\ref{tab:rxn-ole}.

\begin{table}[htbp]
\centering
\caption{Reactions at the m/bl interface ($x = L_{\mathrm{bl}}$).}
\label{tab:rxn-mbl}
\resizebox{\textwidth}{!}{%
\begin{tabular}{clll}
\toprule
\# & Reaction & Role & Lattice \\
\midrule
1 & $m + V_M^{\chi\prime} \rightarrow M_M + v_m + \chi e'$ &
annihilates cation vacancy; cation feed for ol (vacancy path) & conservative \\
2 & $m \rightarrow M_i^{\chi+} + v_m + \chi e'$ &
generates cation interstitial; cation feed for ol (interstitial path) & conservative \\
3 & $m \rightarrow M_M + (\chi/2)\, V_O^{\bullet\bullet} + \chi e'$ &
generates oxygen vacancies: \textbf{bl growth into the metal} & \textbf{non-conservative} \\
\bottomrule
\end{tabular}}
\end{table}

\begin{table}[htbp]
\centering
\caption{Reactions at the bl/ol interface ($x = 0$).}
\label{tab:rxn-blol}
\resizebox{\textwidth}{!}{%
\begin{tabular}{clll}
\toprule
\# & Reaction & Role & Lattice \\
\midrule
4 & $M_M + (\delta/4)\,\mathrm{O_2} + \chi e' \rightarrow V_M^{\chi\prime} + \mathrm{MO}_{\delta/2}$ &
generates cation vacancy; deposits ol (direct $\mathrm{O_2}$ route) & conservative \\
$4'$ & $M_M \rightarrow V_M^{\chi\prime} + M^{\delta+} + (\delta-\chi)\, e'$ &
\begin{tabular}[t]{@{}l@{}}generates cation vacancy; cation ejected\\
(dissolution/re-precipitation route; favored in condensed water)\end{tabular} & conservative \\
5 & $M_i^{\chi+} + (\delta/4)\,\mathrm{O_2} + \chi e' \rightarrow \mathrm{MO}_{\delta/2}$ &
annihilates interstitial; deposits ol & conservative \\
$5'$ & $M_i^{\chi+} \rightarrow M^{\delta+} + (\delta-\chi)\, e'$ &
annihilates interstitial; cation ejected & conservative \\
6 & $V_O^{\bullet\bullet} + \tfrac12 \mathrm{O_2} + 2e' \rightarrow O_O$ &
annihilates oxygen vacancy ($\mathrm{O_2}$ present) & conservative \\
7 & $V_O^{\bullet\bullet} + \mathrm{H_2O} + 2e' \rightarrow O_O + \mathrm{H_2}$ &
annihilates oxygen vacancy (water route) & conservative \\
8 & $V_O^{\bullet\bullet} + \mathrm{H_2O} \rightarrow O_O + 2\mathrm{H^+}$ &
\begin{tabular}[t]{@{}l@{}}annihilates oxygen vacancy, proton release\\
(classical PDM; $\mathrm{H^+}$ available in condensed water)\end{tabular} & conservative \\
9 & $\mathrm{H_2O} + 2e' \rightarrow \tfrac12 \mathrm{H_2} + \mathrm{OH^-}$ &
cathodic partial reaction & (no defect) \\
10 & $\mathrm{MO}_{\chi/2} + \frac{\delta-\chi}{4}\,\mathrm{O_2} \rightarrow \mathrm{MO}_{\delta/2}$ &
converts bl to ol: \textbf{bl destruction} (chemical) & \textbf{non-conservative} \\
$10'$ & $\mathrm{MO}_{\chi/2} + \chi \mathrm{H^+} \rightarrow M^{\delta+} + (\chi/2)\,\mathrm{H_2O} + (\delta-\chi)\, e'$ &
proton-assisted \textbf{bl dissolution} & \textbf{non-conservative} \\
\bottomrule
\end{tabular}}
\end{table}

\begin{table}[htbp]
\centering
\caption{Reaction at the ol/water interface ($x = -L_{\mathrm{ol}}$).}
\label{tab:rxn-ole}
\resizebox{\textwidth}{!}{%
\begin{tabular}{clll}
\toprule
\# & Reaction & Role & Lattice \\
\midrule
11 & $\mathrm{MO}_{\delta/2} + \frac{\theta-\delta}{4}\,\mathrm{O_2} \rightarrow \mathrm{MO}_{\theta/2}(d)$ &
\begin{tabular}[t]{@{}l@{}}\textbf{ol destruction}: dissolution / further oxidation\\
(magnetite $\rightarrow$ hematite, seen in the XPS of Veile et al.)\end{tabular} & \textbf{non-conservative} \\
\bottomrule
\end{tabular}}
\end{table}

Only Reaction~3 (growth) and Reactions~10/10$'$ (destruction) move the
barrier-layer interfaces; only Reaction~11 destroys the outer layer. These four
reactions define the thickness ordinary differential equations.

\subsection{Potential distribution}
\label{sec:potential}

The postulates of the supercritical-water formulation (its Section~2.1) are:
potential drops exist at the m/bl and bl/e interfaces; the field strength
$\varepsilon_f$ in the barrier layer is constant (buffered, independent of $V$
and $L_{\mathrm{bl}}$); and the bl/e drop is linear in $V$ and pH and
independent of $L_{\mathrm{bl}}$:
\begin{align}
\phi_{\mathrm{bl/e}} &= \alpha V + \beta\,\mathrm{pH} + \phi^0 ,
&&\text{(Eq.~(1) of the supercritical-water formulation)} \\
V &= \phi_{\mathrm{m/bl}} + \varepsilon_f L_{\mathrm{bl}} + \phi_{\mathrm{bl/e}} ,
&&\text{(Eq.~(2) of the supercritical-water formulation)} \\
\phi_{\mathrm{m/bl}} &= (1-\alpha) V - \varepsilon_f L_{\mathrm{bl}}
- \beta\,\mathrm{pH} - \phi^0 .
&&\text{(Eq.~(3) of the supercritical-water formulation)}
\end{align}
Here $\alpha$ is the polarizability of the bl/e interface
($\partial \phi_{\mathrm{bl/e}} / \partial V$);
$\beta = \partial \phi_{\mathrm{bl/e}} / \partial\,\mathrm{pH}$
($< 0$, typically $-0.005$ to $-0.05$~V); and $\phi^0$ is the drop at $V = 0$,
$\mathrm{pH} = 0$, $L_{\mathrm{bl}} = 0$. The porous outer layer carries no
potential drop (the cathodic reactions occur at the bl/ol interface via
$\mathrm{O_2}/\mathrm{H_2O}$ transport through pores; there is no remote
cathode current in an autoclave at open circuit).

\subsection{Rate constants (activated-complex theory with partial charges)}
\label{sec:rates}

Every interfacial rate constant takes the form
(Eq.~(8) of the supercritical-water formulation, with the derivation in
Appendix~A of Li et al.\ (2020)):
\begin{equation}
k_i = k_i^{0}\, \exp(a_i V)\, \exp(b_i L_{\mathrm{bl}})\, \exp(c_i\,\mathrm{pH}) ,
\label{eq:rateconstant}
\end{equation}
with the temperature dependence carried by the base standard rate constant
(Eq.~(9) of the supercritical-water formulation):
\begin{equation}
k_i^{00} = k_i^{00\prime}\,
\exp\!\left[ -\frac{\alpha_i\, \Delta G^0_{R_i}}{R}
\left( \frac{1}{T} - \frac{1}{T_0} \right) \right] .
\end{equation}
The coefficients (Table~1 of the supercritical-water formulation), with
$\alpha_i$ the transfer coefficient of Reaction~$i$, are given in
Table~\ref{tab:coeffs}.

\begin{table}[htbp]
\centering
\caption{Exponential coefficients of the interfacial rate constants,
Eq.~(\ref{eq:rateconstant}).}
\label{tab:coeffs}
\begin{tabular}{llll}
\toprule
$i$ & $a_i$ (1/V) & $b_i$ (1/cm) & $c_i$ \\
\midrule
1, 2, 3 (m/bl) & $\alpha_i \chi \gamma (1-\alpha)$ &
$\boldsymbol{-\alpha_i \chi \gamma \varepsilon_f}$ &
$-\alpha_i \chi \gamma \beta$ \\
4, 5, 6, 7, 10, 11 & 0 & 0 & 0 \\
$4'$, $5'$ & $\alpha_i \delta \gamma \alpha$ & 0 & $\alpha_i \delta \gamma \beta$ \\
8, 9 & $2 \alpha_i \gamma \alpha$ & 0 & $2 \alpha_i \gamma \beta$ \\
$10'$ & $\alpha_{10'} (\delta-\chi) \gamma \alpha$ & 0 &
$\alpha_{10'} (\delta-\chi) \gamma \beta$ \\
\bottomrule
\end{tabular}
\end{table}

The m/bl reactions inherit their $V$-, $L_{\mathrm{bl}}$-, and pH-dependence
from $\phi_{\mathrm{m/bl}}$; the bl/ol electrochemical reactions inherit theirs
from $\phi_{\mathrm{bl/e}}$, hence $b_i = 0$ there. The single most important
structural fact is
\begin{equation}
b_3 = -\alpha_3 \chi \gamma \varepsilon_f < 0 :
\end{equation}
the growth reaction decelerates exponentially with film thickness. This is
what generates direct-logarithmic/pseudo-parabolic kinetics and a finite
steady-state thickness.

The steady-state defect flux balances (rates equal at both interfaces for each
defect; Eqs.~(10)--(12) of the supercritical-water formulation) are
\begin{align}
k_1\, C_{V_M}^{L} &= k_4 C_{\mathrm{O}}^{\,n} + k_{4'} , \\
k_2 &= \left( k_5 C_{\mathrm{O}}^{\,n} + k_{5'} \right) C_{M_i}^{0} , \\
\frac{\chi}{2}\, k_3 &= \left( k_6 C_{\mathrm{O}}^{\,m}
+ k_7 C_{\mathrm{H_2O}}^{\,p} + k_8 C_{\mathrm{H_2O}}^{\,p} \right) C_{V_O}^{0} .
\end{align}
The kinetic orders are: $n$ (Reactions 4 and 5 in $C_{\mathrm{O}}$),
$m$ (Reaction 6 in $C_{\mathrm{O}}$), $p$ (Reactions 7--9 in
$C_{\mathrm{H_2O}}$), $q$ (Reaction 10 in $C_{\mathrm{O}}$),
$r$ (Reaction 11 in $C_{\mathrm{O}}$), and $s$ (Reaction $10'$ in the
$\mathrm{H^+}$ activity ratio). Li et al.\ (2020) assume $n = m = q = r = 1/2$.

\subsection{Thickness evolution equations (the governing ODEs)}
\label{sec:odes}

\subsubsection{Barrier layer}

Only the non-conservative reactions contribute
(Eq.~(16) of the supercritical-water formulation, with porosity correction
$\rho^0_{\mathrm{bl}}/\rho_{\mathrm{bl}}$):
\begin{equation}
\frac{\mathrm{d}L_{\mathrm{bl}}}{\mathrm{d}t}
= \Omega_{\mathrm{bl}} \left[ k_3 - k_{10} C_{\mathrm{O}}^{\,q}
- k_{10'} \left( C_{\mathrm{H^+}}/C^0_{\mathrm{H^+}} \right)^{s} \right]
\frac{\rho^0_{\mathrm{bl}}}{\rho_{\mathrm{bl}}} .
\end{equation}
Substituting
$k_3 = k_3^{0} \exp(a_3 V) \exp(b_3 L_{\mathrm{bl}}) \exp(c_3\,\mathrm{pH})$
at constant $V$ and pH (Eqs.~(21)--(23) of the supercritical-water
formulation) gives the working form
\begin{equation}
\boxed{\;
\frac{\mathrm{d}L_{\mathrm{bl}}}{\mathrm{d}t}
= A_{\mathrm{bl}}\, \exp( b_3 L_{\mathrm{bl}} ) - C_{\mathrm{bl}}
\;}
\label{eq:blode}
\end{equation}
with
\begin{align}
A_{\mathrm{bl}} &= \Omega_{\mathrm{bl}}\, k_3^{0}\, \exp(a_3 V)\,
\exp(c_3\,\mathrm{pH})\, \frac{\rho^0_{\mathrm{bl}}}{\rho_{\mathrm{bl}}} ,
&&\text{(Eq.~(22) of the supercritical-water formulation)}
\label{eq:Abl} \\
C_{\mathrm{bl}} &= \Omega_{\mathrm{bl}} \left[ k_{10} C_{\mathrm{O}}^{\,q}
+ k_{10'} \left( C_{\mathrm{H^+}}/C^0_{\mathrm{H^+}} \right)^{s} \right]
\frac{\rho^0_{\mathrm{bl}}}{\rho_{\mathrm{bl}}} .
&&\text{(Eq.~(23) of the supercritical-water formulation)}
\label{eq:Cbl}
\end{align}
(The printed ``$k_9$'' in Eq.~(23) of the supercritical-water formulation is a
typographical error for $k_{10}$, as required by consistency with its
Eqs.~(16), (17), and (44); see Section~\ref{sec:errata}.)

The \textbf{closed-form solution} for constant $A_{\mathrm{bl}}$, $b_3$,
$C_{\mathrm{bl}}$ and initial condition $L_{\mathrm{bl}}(0) = L_0$
(Eqs.~(24) and (39) of the supercritical-water formulation) is
\begin{equation}
L_{\mathrm{bl}}(t) = L_0 - \frac{1}{b_3}
\ln\!\left[ 1 + \frac{A_{\mathrm{bl}}}{C_{\mathrm{bl}}}\, e^{b_3 L_0}
\left( e^{-b_3 C_{\mathrm{bl}} t} - 1 \right) \right]
- C_{\mathrm{bl}}\, t .
\label{eq:blclosed}
\end{equation}

The \textbf{steady state} ($\mathrm{d}L_{\mathrm{bl}}/\mathrm{d}t = 0$;
Eqs.~(17)--(20) of the supercritical-water formulation) is
\begin{equation}
L_{\mathrm{bl,ss}} = \frac{1-\alpha}{\varepsilon_f}\, V
- \frac{\beta}{\varepsilon_f}\, \mathrm{pH}
- \frac{1}{\alpha_3 \chi \varepsilon_f \gamma}
\ln\!\left[ \frac{ k_{10} C_{\mathrm{O}}^{\,q}
+ k_{10'} \left( C_{\mathrm{H^+}}/C^0_{\mathrm{H^+}} \right)^{s} }{ k_3^{0} } \right] .
\end{equation}
$L_{\mathrm{bl,ss}}$ is linear in $V$ and pH --- the PDM signature. In
ultrapure hydrothermal water (0.055~$\mu$S/cm, $\mathrm{pH} \approx 5.6$,
$C_{\mathrm{O}} \approx 10^{-5}$) the destruction terms are very small:
$C_{\mathrm{bl}} \to 0$, the steady state moves to very large $L$, and over a
finite window
\begin{equation}
L_{\mathrm{bl}}(t) \approx L_0 - \frac{1}{b_3}
\ln\!\left( 1 - b_3 A_{\mathrm{bl}}\, e^{b_3 L_0}\, t \right)
\end{equation}
--- direct logarithmic growth, which over 72--480~h is numerically
indistinguishable from the parabolic $t^{0.5}$ law fitted by Veile et al.\
(their Cr layer: $n = 0.4964$, no saturation at 480~h). Quantifying this
equivalence is one of the results of the study.

The limiting-case map (what the data can select between) is:
\begin{itemize}
\item $C_{\mathrm{bl}} = 0$, $|b_3| L \ll 1$: linear growth,
$L \approx L_0 + A_{\mathrm{bl}} t$.
\item $C_{\mathrm{bl}} = 0$, $|b_3| L = O(1)$: direct-logarithmic
(pseudo-parabolic over a decade of $t$).
\item $C_{\mathrm{bl}} > 0$: logarithmic $\to$ saturation at
$L_{\mathrm{bl,ss}}$.
\end{itemize}

\subsubsection{Outer layer and the constant-volume constraint}

The outer layer grows from cations transmitted through the barrier layer
(fluxes $k_1 C_{V_M}^{L}$ via vacancies and $k_2$ via interstitials, deposited
by Reactions 4/4$'$/5/5$'$), plus bl$\to$ol conversion (Reactions 10/10$'$),
minus dissolution (Reaction 11)
(Eq.~(26) of the supercritical-water formulation):
\begin{equation}
\frac{\mathrm{d}L_{\mathrm{ol}}}{\mathrm{d}t}
= \Omega_{\mathrm{ol}} \left[ k_1 C_{V_M}^{L} + k_2
+ k_{10} C_{\mathrm{O}}^{\,q}
+ k_{10'} \left( C_{\mathrm{H^+}}/C^0_{\mathrm{H^+}} \right)^{s}
- k_{11} C_{\mathrm{O}}^{\,r} \right]
\frac{\rho^0_{\mathrm{ol}}}{\rho_{\mathrm{ol}}} .
\label{eq:olode-full}
\end{equation}

\textbf{Constant-volume growth} is imposed (marker experiments show that the
bl/ol interface stays at the original metal surface): the total metal
consumption rate equals the rate of barrier-layer advance into the metal
(Eqs.~(29)--(32) of the supercritical-water formulation):
\begin{align}
\frac{\mathrm{d}L_m}{\mathrm{d}t} &= \Omega_m \left( k_1 C_{V_M}^{L} + k_2 + k_3 \right)
&&\text{(metal destruction, Reactions 1+2+3)} , \\
\frac{\mathrm{d}L_{\mathrm{bl}}}{\mathrm{d}t} &= \Omega_{\mathrm{bl}} \left( k_3 - k_d \right) ,
&&k_d := k_{10} C_{\mathrm{O}}^{\,q}
+ k_{10'} \left( C_{\mathrm{H^+}}/C^0_{\mathrm{H^+}} \right)^{s} ,
\end{align}
and setting $\mathrm{d}L_m/\mathrm{d}t = \mathrm{d}L_{\mathrm{bl}}/\mathrm{d}t$
gives
\begin{equation}
\boxed{\;
k_1 C_{V_M}^{L} + k_2 + k_3 = \mathrm{PBR} \left( k_3 - k_d \right)
\;}
\qquad
\text{(Eq.~(32) of the supercritical-water formulation)} .
\label{eq:cv}
\end{equation}
For an n-type barrier layer (interstitials/oxygen vacancies dominate,
$k_1 C_{V_M}^{L} \ll k_2$) and $k_d \approx 0$
(Eq.~(33) of the supercritical-water formulation):
\begin{equation}
k_2 = (\mathrm{PBR} - 1)\, k_3 .
\end{equation}

Substituting Eq.~(\ref{eq:cv}) into the outer-layer equation
(\ref{eq:olode-full}), the $k_d$ terms combine exactly,
$(\mathrm{PBR}-1) k_3 - \mathrm{PBR}\, k_d + k_d
= (\mathrm{PBR}-1)(k_3 - k_d)$, giving
(Eq.~(41) of the supercritical-water formulation)
\begin{equation}
\boxed{\;
\begin{aligned}
\frac{\mathrm{d}L_{\mathrm{ol}}}{\mathrm{d}t}
&= \Omega_{\mathrm{ol}} \left[ (\mathrm{PBR}-1)(k_3 - k_d)
- k_{11} C_{\mathrm{O}}^{\,r} \right] \\
&= (\mathrm{PBR}-1) \frac{\Omega_{\mathrm{ol}}}{\Omega_{\mathrm{bl}}}
\frac{\mathrm{d}L_{\mathrm{bl}}}{\mathrm{d}t}
- \Omega_{\mathrm{ol}}\, k_{11} C_{\mathrm{O}}^{\,r}\,
\frac{\rho^0_{\mathrm{ol}}}{\rho_{\mathrm{ol}}}
\end{aligned}
\;}
\label{eq:olode}
\end{equation}
so that with negligible barrier-layer destruction the thickness ratio is fixed
by stoichiometry alone (Eqs.~(36) and (37) of the supercritical-water
formulation):
\begin{equation}
\frac{\mathrm{d}L_{\mathrm{ol}}}{\mathrm{d}L_{\mathrm{bl}}}
= (\mathrm{PBR}-1)\, \frac{\Omega_{\mathrm{ol}}}{\Omega_{\mathrm{bl}}}
=: \mathrm{PBR}_{\mathrm{eff}} .
\end{equation}

\textbf{Closed form.} Li et al.\ (2020) write it
(their Eqs.~(42)--(45)) with
$A_{\mathrm{ol}} = (\mathrm{PBR}-1)
(\Omega_{\mathrm{ol}}/\Omega_{\mathrm{bl}}) A_{\mathrm{bl}}$ and
$C_{\mathrm{ol}} = \Omega_{\mathrm{ol}}
[ (\mathrm{PBR}-1) k_d + k_{11} C_{\mathrm{O}}^{\,r} ]$
(so that $\mathrm{d}L_{\mathrm{ol}}/\mathrm{d}t
= A_{\mathrm{ol}} \exp(b_3 L_{\mathrm{bl}}) - C_{\mathrm{ol}}$):
\begin{align}
L_{\mathrm{ol}}(t) &= \frac{A_{\mathrm{ol}}}{A_{\mathrm{bl}} b_3}
\ln\!\left[ A_{\mathrm{bl}}\, e^{b_3 L_{\mathrm{bl}}(t)} - C_{\mathrm{bl}} \right]
- C_{\mathrm{ol}}\, t + D_{\mathrm{ol}} ,
\label{eq:olclosed} \\
D_{\mathrm{ol}} &= -\frac{A_{\mathrm{ol}}}{A_{\mathrm{bl}} b_3}
\ln\!\left[ A_{\mathrm{bl}}\, e^{b_3 L_0} - C_{\mathrm{bl}} \right]
+ L_{\mathrm{ol}}(0) .
\end{align}
Because $\mathrm{d}L_{\mathrm{ol}}/\mathrm{d}t$ is an exact affine function of
$\mathrm{d}L_{\mathrm{bl}}/\mathrm{d}t$, the reduced system also integrates
trivially:
\begin{gather}
L_{\mathrm{ol}}(t) = \mathrm{PBR}_{\mathrm{eff}}
\left( L_{\mathrm{bl}}(t) - L_0 \right) + C_x\, t + L_{\mathrm{ol}}(0) ,
\label{eq:olaffine} \\
\mathrm{PBR}_{\mathrm{eff}} := (\mathrm{PBR}-1)
\frac{\Omega_{\mathrm{ol}}}{\Omega_{\mathrm{bl}}} ,
\qquad
C_x := -\Omega_{\mathrm{ol}}\, k_{11} C_{\mathrm{O}}^{\,r} \;(\le 0) \notag
\end{gather}
--- an independent verification target alongside Eq.~(\ref{eq:olclosed}).

The total scale thickness is
$L_{\mathrm{tot}} = L_{\mathrm{bl}} + L_{\mathrm{ol}}$. The weight gain (not
measurable for the cube specimens of Veile et al., kept for completeness;
Eq.~(28) of the supercritical-water formulation) is
$\mathrm{d}(\Delta w)/\mathrm{d}t
= r_{\mathrm{bl}} \rho_{\mathrm{bl}}\, \mathrm{d}L_{\mathrm{bl}}/\mathrm{d}t
+ r_{\mathrm{ol}} \rho_{\mathrm{ol}}\, \mathrm{d}L_{\mathrm{ol}}/\mathrm{d}t$,
with $r$ the oxygen mass fraction of each oxide; a closed seven-parameter form
is given in Eqs.~(48)--(56) of the supercritical-water formulation.

\subsubsection{Interfacial current (for extensions with explicit corrosion potential)}

The partial currents (Appendix~B of Li et al.\ (2020)) are
$i_1 = \chi F k_1 C_{V_M}^{L}$, $i_2 = \chi F k_2$, $i_3 = \chi F k_3$,
$i_4 = -\chi F k_4 C_{\mathrm{O}}^{\,n}$,
$i_5 = -\chi F k_5 C_{\mathrm{O}}^{\,n} C_{M_i}^{0}$,
$i_6 = -2 F k_6 C_{\mathrm{O}}^{\,m} C_{V_O}^{0}$,
$i_7 = -2 F k_7 C_{\mathrm{H_2O}}^{\,p} C_{V_O}^{0}$,
$i_9 = -2 F k_9 C_{\mathrm{H_2O}}^{\,p}$,
$i_{4'} = (\delta-\chi) F k_{4'}$, etc.
At open circuit the sum of the partial currents vanishes
(Eq.~(46) of the supercritical-water formulation), which fixes
$V = V_{\mathrm{ocp}}$ given the $k_i^{0}$.

\subsection{Reduction to the identifiable apparent-parameter model}
\label{sec:reduction}

At open circuit, $V$ and pH are constant over an exposure, so
$\exp(a_3 V + c_3\,\mathrm{pH})$ is a constant multiplier: it cannot be
separated from $k_3^{0}$ by thickness-versus-time data alone. The observable
model is therefore the \textbf{apparent-parameter system} (identical in
structure to Table~5 of Li et al.\ (2020)):
\begin{align}
\frac{\mathrm{d}L_{\mathrm{bl}}}{\mathrm{d}t}
&= A_{\mathrm{bl}}\, \exp(b_3 L_{\mathrm{bl}}) - C_{\mathrm{bl}} ,
& L_{\mathrm{bl}}(0) &= L_0 ,
\label{eq:app-bl} \\
\frac{\mathrm{d}L_{\mathrm{ol}}}{\mathrm{d}t}
&= \mathrm{PBR}_{\mathrm{eff}}
\left( A_{\mathrm{bl}}\, \exp(b_3 L_{\mathrm{bl}}) - C_{\mathrm{bl}} \right) + C_x ,
& L_{\mathrm{ol}}(0) &= 0 ,
\label{eq:app-ol}
\end{align}
with $C_x := -\Omega_{\mathrm{ol}} k_{11} C_{\mathrm{O}}^{\,r} \le 0$
(outer-layer dissolution / further oxidation only).

The parameter vector is
\begin{multline}
\theta = \big\{
A_{\mathrm{bl}}~\mathrm{[nm/h]},\;
b_3~\mathrm{[1/nm]}\ (<0),\;
C_{\mathrm{bl}}~\mathrm{[nm/h]}\ (\ge 0),\\
\mathrm{PBR}_{\mathrm{eff}}~\mathrm{[-]}\ (>0),\;
C_x~\mathrm{[nm/h]},\;
L_0~\mathrm{[nm]}\ (>0)
\big\}
\end{multline}
--- 6 parameters, 2 ODEs, data at 3 times for 2 layers (plus a Ni consistency
check).

Physical quantities are recovered after fitting via
\begin{align}
\varepsilon_f &= \frac{-b_3}{\alpha_3 \chi \gamma}
&&\text{[from } b_3 = -\alpha_3 \chi \gamma \varepsilon_f \text{]} , \\
k_3^{0} &= \frac{A_{\mathrm{bl}}}{\Omega_{\mathrm{bl}}\, \exp(a_3 V + c_3\,\mathrm{pH})}
&&\text{[order of magnitude, given } \alpha_3, V \text{]} , \\
\mathrm{PBR} &= 1 + \mathrm{PBR}_{\mathrm{eff}}\,
\frac{\Omega_{\mathrm{bl}}}{\Omega_{\mathrm{ol}}} , \\
k_d &= \frac{C_{\mathrm{bl}}}{\Omega_{\mathrm{bl}}} ,
\end{align}
with $\alpha_3 \in [0.1, 0.5]$ (Li et al.\ (2020) fit $\alpha_3 = 0.12$ for
HCM12A) and $\chi = 8/3$, $\gamma = F/RT = 22.63$~V$^{-1}$ at 513.15~K. Molar
volumes per cation:
$\Omega_{\mathrm{bl}} = 46.7/3 = 15.6$~cm$^3$/mol for $\mathrm{FeCr_2O_4}$
(3 cations);
$\Omega_{\mathrm{ol}} = 45.0/3 = 15.0$~cm$^3$/mol for $\mathrm{Fe_3O_4}$;
$\Omega_m \approx 7.1$~cm$^3$/mol for the fcc stainless steel. These are
broadly consistent, as an independent cross-check rather than a
re-derivation, with the theoretical $\mathrm{PBR} = 2.05$
($\mathrm{FeCr_2O_4}$) and $2.10$ ($\mathrm{Fe_3O_4}$) reported in
Table~3 of Li et al.\ (2020) for the same phases against their own
representative steel molar volume; the fit itself recovers
$\mathrm{PBR}_{\mathrm{eff}} = 1.096$, the apparent
outer-to-barrier growth-rate ratio identified from the thickness data
(not $\Omega_{\mathrm{bl}}/\Omega_m$ directly), related to the physical
PBR through the equation above.

Priors/constraints (3 time points cannot determine 7 parameters freely).
The only prior applied in every fit is the Gaussian log-space prior of
Eq.~(3) of the main text: $\ln L_0 \sim
\mathcal{N}(\ln 2~\mathrm{nm},\, 0.60^2)$, whose $\pm2\sigma$ range is
0.6--6.6~nm, matching the air-formed-film range below.

Earlier versions of this analysis also applied
$\ln \mathrm{PBR}_{\mathrm{eff}} \sim \mathcal{N}(\ln 1.05,\, 0.20^2)$,
on the grounds that 1.05 was the stoichiometric value. It is not: the
constant-volume closure giving 1.05 requires the barrier layer to draw
$\approx 30$~at\% Cr from an 18.8~at\% alloy, and a chromium mass balance
with an $\mathrm{FeCr_2O_4}$ barrier gives 2.09 when all non-Ni iron
reports to the outer layer. That prior has been removed;
$\mathrm{PBR}_{\mathrm{eff}}$ is now determined by the likelihood alone
(1.096 for M4). The apparent gap between 1.096 and 2.09 is examined in
the main text and does not survive: the mass-balance figure is
phase-dependent and spans $[0.52, 3.79]$ over the candidate barrier
spinels, and once an outer-layer loss term is admitted the FeCr$_2$O$_4$
value is reached at no cost in $\chi^2$. The physical ranges below
motivate the remaining prior:
\begin{itemize}
\item $L_0$: air-formed film on high-Cr austenitic stainless steel, 1--5~nm
(Li et al.\ (2020): 1.01~nm for 316L).
\item $\mathrm{PBR}_{\mathrm{eff}}$: no prior; fitted freely. The
chromium mass balance admits $[0.52, 3.79]$ depending on the barrier
phase assignment; the likelihood returns 1.096, inside that range.
\item $C_{\mathrm{bl}}$: $\ge 0$, expected $\approx 0$ (ultrapure water); both
$C_{\mathrm{bl}} = 0$ (five-parameter model) and free $C_{\mathrm{bl}}$ are
tested.
\item $b_3$: strictly $< 0$. $|b_3| \sim \alpha_3 \chi \gamma \varepsilon_f$;
with $\varepsilon_f \in [10^2, 10^6]$~V/cm this spans
$|b_3| \in [\sim 7 \times 10^{-5}, \sim 0.7]$~nm$^{-1}$ --- the data must
localize it.
\item $C_x$: free sign, small.
\end{itemize}

The units convention for the implementation is thickness in nm and time in h.
Conversions from the cgs units (cm, s) of Li et al.\ (2020):
1~cm/s $= 3.6 \times 10^{10}$~nm/h; 1~cm$^{-1}$ $= 10^{-7}$~nm$^{-1}$.
Sign conventions respected throughout: $b_3 < 0$, $C_{\mathrm{bl}} \ge 0$,
$\beta < 0$, $D_{\mathrm{ol}} < 0$.

\subsection{Nondimensionalization for the neural spectral-element solver}
\label{sec:nondim}

The characteristic scales are $t_c = 480$~h (the exposure window) and
$L_c = 100$~nm (the observed Cr-layer scale). With $\tau = t/t_c$ and
$\lambda = L/L_c$:
\begin{align}
\frac{\mathrm{d}\lambda_{\mathrm{bl}}}{\mathrm{d}\tau}
&= \hat{A}\, \exp( \hat{b}\, \lambda_{\mathrm{bl}} ) - \hat{C} ,
& \lambda_{\mathrm{bl}}(0) &= \lambda_0 , \\
\frac{\mathrm{d}\lambda_{\mathrm{ol}}}{\mathrm{d}\tau}
&= \mathrm{PBR}_{\mathrm{eff}}
\left( \hat{A}\, \exp( \hat{b}\, \lambda_{\mathrm{bl}} ) - \hat{C} \right)
+ \hat{C}_x ,
\end{align}
with
\begin{equation}
\hat{A} = \frac{A_{\mathrm{bl}}\, t_c}{L_c} , \qquad
\hat{b} = b_3 L_c , \qquad
\hat{C} = \frac{C_{\mathrm{bl}}\, t_c}{L_c} , \qquad
\hat{C}_x = \frac{C_x\, t_c}{L_c} , \qquad
\lambda_0 = \frac{L_0}{L_c} .
\end{equation}
For the expected regime ($L_{\mathrm{bl}}$: 37 $\to$ 134~nm over 72 $\to$
480~h with near-parabolic shape) these groups are $O(0.1$--$10)$: e.g.\ a
pure-logarithmic fit gives $\hat{b} \approx -2 \ldots -4$,
$\hat{A} \approx O(10)$, $\hat{C} \approx 0$. All learnable parameters are
$O(1)$ --- a badly scaled parameter is an unlearnable parameter, so this
scaling is applied from the start.

The residuals are enforced in Crank--Nicolson integral form on the
Gauss--Lobatto--Legendre time grid (otherwise the conditioning of the spectral
first-derivative matrix is the dominant error floor):
\begin{equation}
R_j = \lambda_j - \lambda_{j-1}
- \frac{\Delta t_j}{2} \left( f(\lambda_{j-1}) + f(\lambda_j) \right) ,
\qquad j = 1, \ldots, N_t - 1 ,
\end{equation}
plus the initial-condition residuals $\lambda_{\mathrm{bl}}(0) - \lambda_0$
and $\lambda_{\mathrm{ol}}(0)$.

\subsection{Verification targets}
\label{sec:vtargets}

Four verification targets follow from the derivation (results in
Sections~\ref{sec:baseline} and \ref{sec:qs} below):
\begin{enumerate}
\item \textbf{Closed-form parity}: Radau integration of the apparent-parameter
system of Section~\ref{sec:reduction} must match
Eqs.~(\ref{eq:blclosed}) and (\ref{eq:olclosed})--(\ref{eq:olaffine}) to at
least 10 significant digits for generic parameter sets
$(A_{\mathrm{bl}}, b_3, C_{\mathrm{bl}},$
$\mathrm{PBR}_{\mathrm{eff}}, C_x, L_0)$.
\item \textbf{Regression against Table~5 of Li et al.\ (2020)}: with the HCM12A
values ($L_0 = 4.98 \times 10^{-5}$~cm, $b_3 = -6.08 \times 10^{2}$~cm$^{-1}$,
$A_{\mathrm{bl}} = 3.15 \times 10^{-10}$~cm/s,
$C_{\mathrm{bl}} = 3.66 \times 10^{-12}$~cm/s,
$A_{\mathrm{ol}} = 3.66 \times 10^{-10}$~cm/s,
$C_{\mathrm{ol}} = 4.25 \times 10^{-12}$~cm/s,
$D_{\mathrm{ol}} = -4.18 \times 10^{-2}$) and the corresponding 316L values,
reproduce the published thickness-versus-time behavior.
\item \textbf{Power-law refit}: least squares of $L = k t^n$ on the digitized
data of Veile et al.\ (2024) must recover Cr ($k = 6.521$~nm, $n = 0.4964$)
and Fe ($k = 64.51$~nm, $n = 0.2209$) within digitization error.
\item \textbf{Ni consistency}: the fitted model implies no Ni-layer dynamics;
the data standard-deviation band (31--43~nm, approximately constant) must
contain the assumed-constant line.
\end{enumerate}

\subsection{Known errata in the sources}
\label{sec:errata}

Documented here so they are not propagated:
\begin{itemize}
\item Li et al.\ (2020), Eq.~(23): prints $k_9$ where $k_{10}$ is required
(consistency with Eqs.~(16), (17), and (44) of that work).
\item Veile et al.\ (2024): the exposure times are 72/168/480~h (``78~h'' and
``178~h'' in the Results are typographical errors); the legend of their Fig.~9
swaps the $n$ values between Fe and Cr --- the body text is authoritative
(Fe: $n = 0.2209$; Cr: $n = 0.4964$).
\end{itemize}

\section{Spatial defect-transport model}
\label{sec:spatial-derivation}

This section gives the complete governing equations, boundary conditions,
and initial condition for the spatial extension of Section~2.5 of the main
text, which resolves oxygen- and cation-vacancy transport across the
barrier layer under the field strength identified in
Section~\ref{sec:derivation}. Nothing here is fit to the spatial data;
every kinetic quantity entering this section (the growth prefactor, field
exponent, and barrier-layer thickness $L(t)$) is frozen at its
zero-dimensional value.

\subsection{Species, flux, and nondimensionalization}

Two defect species are resolved on the barrier layer, $x \in [0, L]$ with
the origin at the metal/barrier-layer interface: oxygen vacancies
$c_\text{OV}(x,t)$ (charge $z = +2$, generated at the metal/barrier-layer
interface by the growth reaction and annihilated at the
barrier/outer-layer interface, migrating outward) and cation vacancies
$c_\text{CV}(x,t)$ (charge $z = -\chi$ with $\chi = 8/3$ the
spinel-averaged cation valence, generated at the barrier/outer-layer
interface and annihilated at the metal/barrier-layer interface, migrating
inward). Each species obeys a Nernst--Planck flux under the constant
field $\varepsilon_f$ identified in the zero-dimensional fit, with
Einstein mobility $u = D\gamma$, $\gamma = F/(RT)$:
\begin{equation}
J = -D\,\frac{\partial c}{\partial x} + z\gamma\varepsilon_f D\,c.
\label{eq:nernstplanck}
\end{equation}
Nondimensionalizing depth as $\xi = x/L$, concentration as
$\hat{c} = c/c_0$ with the migration scale
$c_0 = |J|/|v|$, $v = z\gamma\varepsilon_f D$ (the migration velocity),
and flux as $\hat{J} = J/(|v|c_0) = \operatorname{sign}(v)$, the steady
equation becomes
\begin{equation}
\frac{d}{d\xi}\!\left[-\frac{1}{\mathrm{Pe}}\frac{d\hat{c}}{d\xi} + s\hat{c}\right] = 0,
\qquad
\mathrm{Pe} = |z|\gamma\varepsilon_f L, \qquad s = \operatorname{sign}(z).
\label{eq:steadypde}
\end{equation}
As discussed in the main text, the diffusivity $D$ cancels from the
P\'eclet number $\mathrm{Pe}$: the nondimensional steady profile shape
depends only on quantities the zero-dimensional fit already fixes, and
$D$ enters only the dimensional concentration scale $c_0$ (invisible to
a composition measurement such as EDX).

\subsection{Boundary conditions}

Both species obey a Robin (interfacial-kinetics) boundary condition at
their annihilation interface, with $\kappa$ the ratio of the interfacial
rate constant to the migration velocity:
\begin{align}
\text{OV (annihilation at } \xi=1\text{):}\quad & J = k_a\,c(L) \;\Longrightarrow\; \hat{c}(1) = 1/\kappa_a,
\label{eq:bc-ov}\\
\text{CV (annihilation at } \xi=0\text{):}\quad & |J| = k_1\,c(0) \;\Longrightarrow\; \hat{c}(0) = 1/\kappa_1.
\label{eq:bc-cv}
\end{align}
The generation flux at the opposite interface is anchored to the
zero-dimensional growth kinetics rather than independently fit: the
oxygen-vacancy generation flux is $J_\text{OV} = (\chi/2)\,k_3(L)$ with
$k_3(L)$ the barrier-layer growth rate constant at thickness $L$
(mol\,cm$^{-2}$\,s$^{-1}$, converted from the fitted $dL_\text{bl}/dt$),
and the cation-vacancy generation flux is anchored to it by an assumed
ratio, $|J_\text{CV}| = r_\text{CV} J_\text{OV}$.

The parameters $\kappa_a$, $\kappa_1$, $r_\text{CV}$, and the two
diffusivities $D_\text{OV}$, $D_\text{CV}$ are \emph{assumed}, not
identified: per the synthetic identifiability gate of Section~2.5 of the
main text, three exposure times cannot constrain interfacial transport
rate constants, so these are treated as forward-only inputs, documented
here for reproducibility rather than reported as findings. The values
used throughout are $\varepsilon_f = 1.73\times10^4$~V/cm (the
zero-dimensional fit), $\Omega_\text{bl} = 15.6$~cm$^3$/mol,
$D_\text{OV} = 10^{-15}$~cm$^2$/s, $D_\text{CV} = 10^{-16}$~cm$^2$/s,
$\kappa_a = \kappa_1 = 1.25$, and $r_\text{CV} = 0.5$.

\subsection{Analytic steady solution}

Because the nondimensional flux is constant in steady state
($\hat{J} = s$ everywhere), Eq.~\eqref{eq:steadypde} admits an exact
closed-form solution: a uniform plateau at $\hat{c} = 1$ (the particular
solution) plus a homogeneous exponential mode that satisfies the Robin
condition at the annihilation boundary $\xi_\text{bc} \in \{0, 1\}$:
\begin{equation}
\hat{c}(\xi) = 1 + \left(\hat{c}_\text{bc} - 1\right)
\exp\!\big[s\,\mathrm{Pe}\,(\xi - \xi_\text{bc})\big].
\label{eq:steadyanalytic}
\end{equation}
This is the reference solution the Newton and NSEM steady solves are
verified against (Section~\ref{sec:baseline}).

\subsection{Transient (moving-boundary) problem}

The transient problem is solved on a Landau-transformed fixed domain
$\xi \in [0,1]$ that tracks the growing barrier-layer thickness $L(t)$,
with $\tau = t/t_c$ the nondimensional time ($t_c = 480$~h, the exposure
window). Writing $\mathrm{Pe}(\tau) = |z|\gamma\varepsilon_f L(\tau)$,
$\hat{\jmath}(\tau) = k_3(L(\tau))/k_3(L_\text{ref})$ the growth-flux
ratio to a reference thickness $L_\text{ref} = L(\tau{=}1)$,
$\mathrm{adv}(\tau) = d\ln L/d\tau$ the Landau advection strength, and
$\mathrm{mig}(\tau) = t_c/\tau_\text{mig}(L(\tau))$ the flux-divergence
prefactor ($\tau_\text{mig} = L/|v|$ the migration transit time), the
governing equation is
\begin{equation}
\frac{\partial \hat{c}}{\partial \tau}
= \xi\,\mathrm{adv}(\tau)\,\frac{\partial \hat{c}}{\partial \xi}
- \mathrm{mig}(\tau)\,\frac{\partial \hat{J}}{\partial \xi},
\qquad
\hat{J} = -\frac{1}{\mathrm{Pe}(\tau)}\frac{\partial \hat{c}}{\partial \xi} + s\hat{c}.
\label{eq:transientpde}
\end{equation}
\textbf{Boundary conditions} (per time node): at the generation
boundary, $\hat{J} = s\,\hat{\jmath}(\tau)$ (the zero-dimensional
growth-flux anchor); at the annihilation boundary,
$\hat{J} = s\,\kappa\,\hat{c}$ (the same Robin closure as the steady
problem). \textbf{Initial condition}: the instantaneous quasi-steady
profile at $\tau_0$, $\hat{c}(\xi, \tau_0) = \hat{\jmath}(\tau_0)\,\hat{c}_\text{ss}(\xi;\mathrm{Pe}(\tau_0))$,
with $\hat{c}_\text{ss}$ the analytic steady shape of
Eq.~\eqref{eq:steadyanalytic}. Quasi-steadiness is quantified in
Section~\ref{sec:qs} by comparing the transient solution at $\tau=1$
against the instantaneous quasi-steady profile at that same moment.

\subsection{Synthetic-EDX instrument and heterogeneity model for the
identifiability gate}
\label{sec:instrument}

The synthetic identifiability gate of the main text (20 replicate
experiments per candidate parameter set) generates each replicate with
the following instrument and heterogeneity model, anchored to the
scan-to-scan scatter of Veile et al.'s replicate line scans:
\begin{itemize}
\item \textbf{Spatial sampling:} composition profiles sampled every
2~nm over $x \in [-120, 220]$~nm (170 points per scan), four element
traces (O, Cr, Fe, Ni) per scan, three scans per replicate experiment.
\item \textbf{EDX spatial resolution:} Gaussian smoothing with
$\sigma = 2/2.355 \approx 0.85$~nm (2~nm full width at half maximum).
\item \textbf{Per-scan position jitter:} each scan is rigidly shifted
by a Gaussian offset with $\sigma = 0.945$~nm.
\item \textbf{Compositional noise:} independent Gaussian noise per
point, standard deviation $\max(0.02\,|v|,\ 0.2)$ weight percent (2\%
relative, floored at 0.2 absolute weight percent).
\item \textbf{Scan-to-scan physical heterogeneity:} each scan draws its
own realization of the shape parameters from log-normal distributions
with per-parameter coefficients of variation $L_\text{bl}$: 0.04,
$w_\text{int}$: 0.30, $w_\text{out}$: 0.30, $A_\text{Ni}$: 0.15,
$w_\text{Ni}$: 0.14, $f_\text{Cr,bl}$: 0.05, $g_\text{bl}$: 0.30,
calibrated to the replicate-scan scatter of Veile et al.'s Fig.~9.
Omitting this heterogeneity term (treating instrument noise as the only
uncertainty source) produces a spuriously optimistic identifiability
verdict for every parameter.
\item \textbf{Verdict thresholds:} a parameter is classed identifiable
if the relative standard deviation of its 20 replicate recoveries is
below 0.25, weakly identifiable below 0.50, and unidentifiable
otherwise. The systematic (accuracy, as opposed to precision) error of
EDX quantification, of order $\pm10\%$ relative, is not simulated; it
floors the absolute-composition parameters $f_\text{Cr,bl}$ and
$A_\text{Ni}$ independently of the statistical verdict.
\end{itemize}

\section{Nickel-exclusion moving-frame model}
\label{sec:ni-model}

The nickel enrichment zone at the receding metal/barrier-layer interface
(Section~2.6 of the main text) is modeled as Wagner-type solute
enrichment ahead of a moving boundary, in the interface-attached frame
$\xi$ (depth into the metal ahead of the interface, $\xi \ge 0$):
\begin{equation}
\frac{\partial c}{\partial t} = D_\text{Ni,eff}\,\frac{\partial^2 c}{\partial \xi^2}
+ v_m(t)\,\frac{\partial c}{\partial \xi},
\label{eq:wagner}
\end{equation}
where the metal streams toward $\xi=0$ at the combined interface
recession speed
\begin{equation}
v_m(t) = \frac{1}{\mathrm{PBR}_\text{bl}}\frac{dL_\text{bl}}{dt}
+ \phi_\text{ol,supply}\,\frac{1}{\mathrm{PBR}_\text{ol}}\frac{dL_\text{ol}}{dt},
\label{eq:vm}
\end{equation}
where $\mathrm{PBR}_\text{bl} := \Omega_\text{bl}/\Omega_m$ and
$\mathrm{PBR}_\text{ol} := \Omega_\text{ol}/\Omega_m$ are the
phase-specific Pilling--Bedworth ratios of the barrier and outer
layers against the metal (distinct from the single unsubscripted
$\mathrm{PBR}$ of Section~\ref{sec:derivation}, which refers to the
barrier layer only). Numerically $\mathrm{PBR}_\text{bl} = 2.05$ and
$\mathrm{PBR}_\text{ol} = 2.10$ (FeCr$_2$O$_4$ and Fe$_3$O$_4$
molar-volume ratios, Li et al.\ 2020 Table~3) and both layer growth
rates are taken from the frozen zero-dimensional kinetics. The free parameter $\phi_\text{ol,supply}$ is
the fraction of the outer layer's cation supply drawn from the matrix
through the barrier layer (as opposed to re-precipitation of
already-dissolved species); $r_\text{Ni}$, the nickel incorporation
ratio, enters through the boundary condition below.

\textbf{Boundary condition} (rejection flux at the moving interface,
$\xi=0$):
\begin{equation}
D_\text{Ni,eff}\,\frac{\partial c}{\partial \xi}\bigg|_{\xi=0}
= -v_m(t)\,(1 - r_\text{Ni})\,c(0,t),
\label{eq:ni-bc}
\end{equation}
which in the steady limit reduces to the classical rejection result
$c(0) = c_m/r_\text{Ni}$. \textbf{Far-field condition}:
$c(\xi \to \infty, t) = c_m$, the bulk matrix nickel content
($c_m = 10$~wt\%). \textbf{Initial condition}: $c(\xi, 0) = c_m$
everywhere (no pre-existing enrichment).

The model is solved by method-of-lines finite differencing on
$\xi \in [0, 300]$~nm with 150 nodes and Radau time integration
(relative tolerance $10^{-8}$, absolute tolerance $10^{-10}$), not by
the NSEM machinery, since the three free parameters
($D_\text{Ni,eff}$, $r_\text{Ni}$, $\phi_\text{ol,supply}$) are fit
directly to four observables (zone widths at 72/168/480~h and the 72~h
peak amplitude) by nonlinear least squares. Observed features are
extracted from the solved profile using Veile et al.'s own convention:
zone width $W(t)$ is the extent over which $c > c_m + 2$~wt\%, and
amplitude $A(t) = c(0,t)$.

\subsection{Fitted values under the rejected constant-mobility model}
\label{sec:ni-fit}

The constant-mobility model above is rejected by the data:
$\chi^2 = 25.71$ on one degree of freedom, with the discrepancy in the
\emph{shape} of $W(t)$ (predicted monotonically growing, measured
approximately flat). Because parameter estimates and nested-model
$\Delta\chi^2$ comparisons obtained under a structurally rejected model
are not interpretable as estimates or as evidence, the fitted values are
recorded here for completeness and reproducibility only, and are not used
to support any claim in the main text. Table~\ref{tab:ni-fit} gives them.

\begin{table}[htbp]
\centering
\caption{Moving-frame nickel fit to the measured zone widths and 72-h
peak amplitude, \emph{under a model the data reject in shape}. The values
are recorded for reproducibility and carry no inferential weight.
$D_\text{Ni,eff}$ and $r_\text{Ni}$ are jointly, not individually,
constrained: both are pinned in combination (large $\Delta\chi^2$ along
most directions) but trade off along one correlated direction, so
neither has a one-dimensional profile interval.}
\label{tab:ni-fit}
\begin{tabular}{lcccl}
\toprule
Parameter & Fit & Profile 1$\sigma$ & Max.\ $\Delta\chi^2$ & Note \\
\midrule
$D_\text{Ni,eff}$ (nm$^2$/h) & 3.905 & correlated with $r_\text{Ni}$ & 101.6 & jointly constrained \\
$r_\text{Ni}$ & 0.2954 & correlated with $D_\text{Ni,eff}$ & 123.6 & jointly constrained \\
$\phi_\text{ol,supply}$ & 0.5708 & [0.579, 0.895] & 3.6 & weak \\
\multicolumn{5}{l}{$\chi^2$/dof $= 25.71/1$; with $\phi_\text{ol,supply}$ fixed to zero: $\chi^2 = 32.1$.} \\
\bottomrule
\end{tabular}
\end{table}

\subsection{Identifiability gate for the time-dependent mobility closures}
\label{sec:ni-closures}

Two one-parameter extensions were tested against the shape failure: a
defect-flux-coupled effective mobility, $D_\text{Ni,eff} \propto
(\mathrm{d}L_\text{bl}/\mathrm{d}t)^{n_\text{flux}}$, tied to the
identified growth kinetics; and a finite-capacity interfacial trapping
closure with saturable trap density $S_\text{cap}$. Each was subjected
to the same twenty-replicate synthetic recovery gate used for the
spatial shape parameters (Section~\ref{sec:instrument}) \emph{before}
being fitted to real data. Both fail (Table~\ref{tab:ni-closures}), so
neither was carried to the real observables.

\begin{table}[htbp]
\centering
\caption{Identifiability gate for the two candidate time-dependent
Ni-mobility closures; both fail on synthetic data at this dataset's
noise level and exposure-time count. RSD is the 20-replicate relative
standard deviation of the recovered parameter.}
\label{tab:ni-closures}
\begin{tabular}{llccl}
\toprule
Closure & New param. & Truth & RSD & Verdict \\
\midrule
Flux-coupled & $n_\text{flux}$ & 1.5 & 27\% & not identifiable (destabilizes $D_\text{Ni}$ to 37\%) \\
Finite-capacity & $S_\text{cap}$ & 1125 & 179\% & not identifiable \\
\bottomrule
\end{tabular}
\end{table}

A forward-only scan of $n_\text{flux}$ over $[-2, 3]$ at the frozen base
kinetics confirms the negative result independently: no single value is
simultaneously shape-correct (flat or non-monotonic $W(t)$) and a
quantitative improvement on the constant-mobility $\chi^2$.

\section{NSEM implementation and training details}
\label{sec:nsem-details}

This section documents every equation, loss term, boundary/initial
condition, and hyperparameter of the four NSEM training configurations
used in this paper, at the level of detail needed to reproduce them
independently of the released code.

\subsection{Shared representation}

Every configuration represents its field(s) on a discrete-variational-
representation (DVR) collocation node set (Gauss--Lobatto--Legendre
nodes, optionally log-clustered toward one end via a compression
parameter $\alpha$) using an element network with hidden width 32 and 3
layers, either a standard multilayer perceptron or a Kolmogorov--Arnold
network backbone (polynomial degree 4). Thickness fields are passed
through a softplus nonlinearity plus a floor of $10^{-6}$ to guarantee
strict positivity. All training uses double-precision (float64)
arithmetic.

\subsection{Configuration 1: zero-dimensional forward solve}
\label{sec:config-forward}

Two time-only element networks, $\lambda_\text{bl}(\tau)$ and
$\lambda_\text{ol}(\tau)$ (nondimensional barrier- and outer-layer
thickness), on $N_t = 24$ DVR nodes over $\tau \in [0, 1]$
(log-clustered, compression $\alpha_t = 1.5$, mapping toward $\tau=0$).
Training proceeds in two stages: (1) an initial-condition pretrain
(800 Adam steps at learning rate $10^{-2}$ plus an L-BFGS polish) that
drives both networks to their constant initial-condition values,
$\lambda_\text{bl} = \lambda_0$ and $\lambda_\text{ol} = \lambda_{\text{ol},0}$
everywhere; (2) physics training against the Crank--Nicolson residuals
of the nondimensional ODE system of Section~\ref{sec:derivation}, with
three loss terms (\texttt{ode\_bl}, \texttt{ode\_ol}, \texttt{ic})
combined by the BRDR aggregator (equal initial weights, adaptively
rebalanced during training) and optimized for 2000 Adam steps
(learning rate $10^{-3}$) followed by 3 L-BFGS phases (max.\ 400
iterations each, strong-Wolfe line search, gradient/step tolerances
$10^{-13}$/$10^{-15}$). The initial-condition loss carries a fixed
multiplier $\lambda_\text{ic} = 10$. Verification: relative
$L^\infty$ error against the closed-form solution, acceptance
$\le 0.5\%$.

\subsection{Configuration 2: zero-dimensional inverse, hard mode}
\label{sec:config-hard}

The kinetic parameters
$\{A_\text{bl}, b_3, \mathrm{PBR}_\text{eff}, L_0, L_{\text{ol},0}\}$
(optionally $C_\text{bl}$) are the only learnable quantities,
parameterized in log space (negative-log for $b_3$) to guarantee sign
and positivity by construction. The thickness trajectories are
generated by differentiable fourth-order Runge--Kutta integration of
the nondimensional ODE (8 sub-steps per grid interval) through these
parameters, then the affine outer-layer relation, so the governing
equations are satisfied exactly and only a data-plus-prior loss is
minimized:
\begin{equation}
\mathcal{L}_\text{hard} =
\ell_\text{data}(\lambda_\text{bl}; t_\text{Cr}, \bar{L}_\text{Cr}, \sigma_\text{Cr})\, n_\text{Cr}
+ \ell_\text{data}(\lambda_\text{ol}; t_\text{Fe}, \bar{L}_\text{Fe}, \sigma_\text{Fe})\, n_\text{Fe}
+ \lambda_\text{prior}\Big(\tfrac{\ln L_0-\ln(2/L_c)}{0.60}\Big)^{\!2},
\label{eq:hardloss}
\end{equation}
where $\ell_\text{data}$ is the mean squared, error-normalized
misfit between the network trajectory (interpolated to the data times)
and the observations, $n_\text{Cr}=n_\text{Fe}=3$ are the exposure-time
counts (matching the main-text $\chi^2$ normalization), and
$\lambda_\text{prior}=1$. Optimized with 4000 Adam steps
(learning rate $10^{-2}$) followed by one L-BFGS phase (max.\ 400
iterations). This is the reference inverse mode for the
zero-dimensional model; the real-data run initializes at $2\times$ the
reference parameter values (\texttt{init\_scale}=2.0) to verify
convergence is not an artifact of a favorable initialization.

\subsection{Configuration 3: zero-dimensional inverse, soft mode}
\label{sec:config-soft}

Identical field representation and pretraining to
Section~\ref{sec:config-forward} (kinetics frozen during pretrain;
pretrain targets are \emph{linear ramps} from the initial condition to
the last data point, not constants, so the kinetic parameters receive a
nonzero time-variation gradient once joint training begins), then joint
training of the field networks and the kinetic parameters against six
loss terms combined by BRDR:
\begin{equation}
\mathcal{L}_\text{soft} =
\lambda_\text{phys}\big(\ell_\text{ode,bl} + \ell_\text{ode,ol}\big)
+ \lambda_\text{ic}\,\ell_\text{ic}
+ \lambda_\text{data}\big(\ell_\text{data,bl} + \ell_\text{data,ol}\big)
+ \lambda_\text{prior}\,\ell_\text{prior},
\label{eq:softloss}
\end{equation}
with $\lambda_\text{phys} = 1000$ (fixed multiplier on the ODE
residual terms; without it the joint optimum leaves residual slack
that is absorbed into the trajectory instead of the kinetics ---
the origin of failure mode F-1), $\lambda_\text{ic} = 10$,
$\lambda_\text{data} = 100$, and $\lambda_\text{prior} = 1$. The two
exhibits reported in the main text (Table~6) use
$\lambda_\text{data}=10$ as an independent reweighting check. Optimized
with 4000 Adam steps (learning rate $10^{-3}$) and 3 L-BFGS phases
(max.\ 400 iterations each). The self-consistency check re-solves the
recovered kinetic parameters through the exact closed-form/RK4 forward
solve and compares the resulting $\chi^2$ to the trajectory $\chi^2$;
consistency is declared satisfied if
$\chi^2_\text{exact} < \max(2\chi^2_\text{traj},\, \chi^2_\text{traj}+2)$.
The constants in this acceptance rule are a working heuristic (tolerate
up to a doubling, or an absolute slack of two $\chi^2$ units for
near-zero trajectory misfits), not a derived statistical threshold; the
failure-mode exhibits in the main text exceed it by more than an order
of magnitude, so the qualitative verdict does not depend on the
specific constants. The specific $\chi^2$ values of a soft-mode run
vary from training run to training run (an earlier run of the
$\lambda_\text{data}=10$ configuration landed at trajectory
$\chi^2 = 0.004$ against closed-form $\chi^2 = 47.3$), but the
orders-of-magnitude inconsistency signature is stable across runs; the
main text quotes single-seed (seed~0) values.

\subsection{Configuration 4: spatial extension (steady and transient)}
\label{sec:config-spatial}

\textbf{Steady} (Section~\ref{sec:spatial-derivation}): a single
field network on $N_x = 24$ DVR nodes, uniformly spaced on
$\xi \in [0,1]$ ($\alpha_x = 0$; the P\'eclet numbers in this problem,
10--14, produce only mild boundary layers, so uniform spacing suffices).
The barrier-layer thickness is fixed at its 480-h value from the
zero-dimensional fit. The loss is the first-integral form of the
steady equation --- $(\hat{J}-s)^2$ averaged over all nodes, which
penalizes flux non-conservation directly rather than differentiating a
second-derivative residual --- plus the squared Robin-boundary
violation at the annihilation node, trained for 2000 Adam steps
(learning rate $10^{-3}$) after a 500-step pretrain to the boundary
value, then L-BFGS polish. The independent Newton reference solves the
same first-integral system on the identical node set by Newton
iteration (converges in one step, since the discrete system is linear;
verified to residual $<10^{-13}$).

\textbf{Transient}: field network on the tensor-product grid of the
same $N_x=24$ spatial DVR nodes and $N_t=16$ temporal DVR nodes
($\alpha_t=1.5$, log-clustered, $\tau \in [\tau_0, 1]$ with
$\tau_0=0.05$). Crank--Nicolson residuals of
Eq.~\eqref{eq:transientpde} are combined with the two per-time-node
boundary residuals and the initial-condition residual (four loss terms:
\texttt{pde}, \texttt{bc\_gen}, \texttt{bc\_ann}, \texttt{ic}), BRDR-
aggregated and optimized identically to the steady case.

\subsection{Hyperparameter summary}

Table~\ref{tab:hyperparams} collects every numerical hyperparameter
referenced above for exact reproduction.

\begin{table}[htbp]
\centering
\caption{NSEM hyperparameters by configuration.}
\label{tab:hyperparams}
\begin{tabular}{lcccc}
\toprule
Parameter & Forward & Inverse (hard) & Inverse (soft) & Spatial (steady/transient) \\
\midrule
$N_t$ (time nodes) & 24 & 24 & 24 & --- / 16 \\
$N_x$ (space nodes) & --- & --- & --- & 24 / 24 \\
$\alpha_t$, $\alpha_x$ & 1.5 & 1.5 & 1.5 & 0.0 / 1.5 \\
Backbone & MLP/KAN & --- & MLP/KAN & MLP \\
Hidden width, layers & 32, 3 & --- & 32, 3 & 32, 3 \\
Polynomial degree (KAN) & 4 & --- & 4 & 4 \\
Pretrain steps & 800 & --- & 800 & 500 \\
Pretrain learning rate & $10^{-2}$ & --- & $10^{-2}$ & $10^{-2}$ \\
Adam steps & 2000 & 4000 & 4000 & 2000 \\
Adam learning rate & $10^{-3}$ & $10^{-2}$ & $10^{-3}$ & $10^{-3}$ \\
L-BFGS phases $\times$ max.\ iter. & 3 $\times$ 400 & 1 $\times$ 400 & 3 $\times$ 400 & 3 $\times$ 400 \\
L-BFGS tolerances (grad, change) & $10^{-13}$, $10^{-15}$ & $10^{-13}$, $10^{-15}$ & $10^{-13}$, $10^{-15}$ & $10^{-13}$, $10^{-15}$ \\
$\lambda_\text{ic}$ & 10 & --- & 10 & 10 \\
$\lambda_\text{data}$ & --- & (built into $\ell_\text{data}$) & 100 (or 10) & --- \\
$\lambda_\text{phys}$ & --- & --- & 1000 & --- \\
$\lambda_\text{prior}$ & --- & 1 & 1 & --- \\
Random seed & 0 & 0 & 0 & 0 \\
\bottomrule
\end{tabular}
\end{table}

\section{Baseline numerical verification}
\label{sec:baseline}

Figure~\ref{fig:baseline} summarizes the baseline verification of the forward
model. The closed-form solutions, Eqs.~(\ref{eq:blclosed}) and
(\ref{eq:olclosed}), agree with a stiff Radau integration of the
apparent-parameter ODE system to better than
$6 \times 10^{-11}$~nm on the demonstration parameter set
(maximum absolute deviation $4.63 \times 10^{-11}$~nm for the barrier layer
and $5.07 \times 10^{-11}$~nm for the outer layer), i.e.\ machine-level parity.
In addition, the implementation reproduces the published
thickness-versus-time behavior when driven by the Table-5 apparent parameters
of Li et al.\ (2020) for HCM12A and 316L at 500~$^\circ$C in supercritical
water, confirming that the adapted model regresses exactly to the
supercritical-water formulation in its original regime.

\begin{figure}[htbp]
\centering
\includegraphics[width=\textwidth]{figures/fig2_baseline_verification.png}
\caption{Baseline numerical verification. Closed-form solutions versus Radau
integration for the barrier and outer layers (parity better than
$6 \times 10^{-11}$~nm), and regression against the Table-5 apparent
parameters of Li et al.\ (2020) for HCM12A and 316L at 500~$^\circ$C in
supercritical water.}
\label{fig:baseline}
\end{figure}

\section{Quasi-steady approximation error}
\label{sec:qs}

The interface-reaction-controlled model of Section~\ref{sec:derivation} treats
the defect concentration profiles inside the barrier layer as quasi-steady.
Figure~\ref{fig:transient} quantifies the error incurred by this assumption
using a full moving-boundary transport solve. At the end of the 480~h exposure
window, the oxygen-vacancy species --- whose migration transit time across the
film is approximately 5~h --- lags its quasi-steady profile shape by only
0.69\%, while the slower cation-vacancy species (transit time approximately
36~h) lags by 5.79\%. The quasi-steady closure is therefore an excellent
approximation for the species that controls film growth, and an acceptable one
for the slowest defect over the experimental window.

\begin{figure}[htbp]
\centering
\includegraphics[width=\textwidth]{figures/fig9_transient_qs.png}
\caption{Quasi-steady approximation error from the moving-boundary transient
solve. At $t = 480$~h the oxygen-vacancy profile (migration transit time
$\approx 5$~h) deviates from its quasi-steady shape by 0.69\%; the
cation-vacancy profile ($\approx 36$~h) by 5.79\%.}
\label{fig:transient}
\end{figure}

\section{Neural-backbone parity}
\label{sec:backbone}

The inverse-identification results in the main text are independent of the
choice of neural backbone. Trained on the forward problem and compared against
the closed-form solutions of Section~\ref{sec:derivation}, the multilayer
perceptron (MLP) backbone achieves relative $L^\infty$ errors of 0.14\% for
the barrier layer and 0.12\% for the outer layer, while the
Kolmogorov--Arnold backbone achieves 0.11\% and 0.04\%, respectively. Both
backbones are well inside the 0.5\% acceptance bar, so the parameter estimates
and uncertainty statements reported in the main text do not depend on the
backbone architecture.
